% === Overview: Hybrid Problem-Evolution Graph ===
% 7 horizontal lanes (one per thread), shared x-axis 2017--2026,
% landmark papers as nodes, dashed cross-thread arrows.

\begin{figure}[htbp]
\centering
\resizebox{\textwidth}{!}{%
\begin{tikzpicture}[
    >=Stealth,
    node distance=0.6cm,
    paper/.style={
        rounded corners=2pt,
        font=\footnotesize\bfseries,
        inner sep=4pt,
        minimum height=0.65cm,
        text=white,
        drop shadow={shadow xshift=0.5pt, shadow yshift=-0.5pt, opacity=0.3},
    },
    crossarrow/.style={
        ->,
        dashed,
        very thick,
        shorten >=2pt,
        shorten <=2pt,
    },
    arrowlabel/.style={
        font=\tiny\itshape,
        fill=white,
        fill opacity=0.85,
        text opacity=1,
        inner sep=1pt,
        rounded corners=1pt,
    },
    lanelabel/.style={
        font=\small\bfseries,
        anchor=east,
        text width=4.2cm,
        align=right,
    },
]

% --- Geometry ---
\def\laneheight{1.7}   % vertical spacing between lanes
\def\xmin{0}
\def\xmax{20}          % total width
% Year mapping: 2017 = 0, 2026 = 20  =>  1 year = 2 units
\newcommand{\yearx}[1]{\fpeval{(#1-2017)*2}}

% --- Year axis ---
\foreach \y in {2017,2018,2019,2020,2021,2022,2023,2024,2025,2026} {
    \draw[gray!40] (\yearx{\y}, 0.5*\laneheight) -- (\yearx{\y}, -6.5*\laneheight);
    \node[font=\footnotesize, gray!80, anchor=south] at (\yearx{\y}, 0.7*\laneheight) {\y};
}

% --- Lane backgrounds ---
\foreach \i/\col in {0/thread1,1/thread2,2/thread3,3/thread4,4/thread5,5/thread6,6/thread7} {
    \fill[\col, opacity=0.06]
        (\xmin, -\i*\laneheight+0.45*\laneheight)
        rectangle
        (\xmax, -\i*\laneheight-0.45*\laneheight);
}

% =====================================================================
% Thread 1: Instruction Following & SFT (lane y=0)
% =====================================================================
\def\laneA{0}
\node[lanelabel, text=thread1] at (-0.3, \laneA*\laneheight)
    {T1: Instruction\\Following \& SFT};

\node[paper, fill=thread1] (flan)
    at (\yearx{2021}+0.5, \laneA*\laneheight) {FLAN};
\node[paper, fill=thread1] (instructgpt-sft)
    at (\yearx{2022}+0.8, \laneA*\laneheight) {InstructGPT};
\node[paper, fill=thread1] (selfinstruct)
    at (\yearx{2023}+1.2, \laneA*\laneheight) {Self-Inst.};
\node[paper, fill=thread1] (lima)
    at (\yearx{2024}+0.5, \laneA*\laneheight) {LIMA};

% =====================================================================
% Thread 2: Reward Modeling & Preference Optimization (lane y=-1)
% =====================================================================
\def\laneB{-1}
\node[lanelabel, text=thread2] at (-0.3, \laneB*\laneheight)
    {T2: Reward Modeling \&\\Preference Opt.};

\node[paper, fill=thread2] (rlhf-og)
    at (\yearx{2017}+1.0, \laneB*\laneheight) {RLHF};
\node[paper, fill=thread2] (instructgpt-rlhf)
    at (\yearx{2022}+0.8, \laneB*\laneheight) {InstructGPT};
\node[paper, fill=thread2] (dpo)
    at (\yearx{2023}+0.6, \laneB*\laneheight) {DPO};
\node[paper, fill=thread2] (grpo)
    at (\yearx{2024}+1.0, \laneB*\laneheight) {GRPO};
\node[paper, fill=thread2] (dapo)
    at (\yearx{2025}+1.0, \laneB*\laneheight) {DAPO/RLVR};

% =====================================================================
% Thread 3: Safety & Alignment (lane y=-2)
% =====================================================================
\def\laneC{-2}
\node[lanelabel, text=thread3] at (-0.3, \laneC*\laneheight)
    {T3: Safety \&\\Alignment};

\node[paper, fill=thread3] (cai)
    at (\yearx{2022}+0.8, \laneC*\laneheight) {Const.\ AI};
\node[paper, fill=thread3] (saferlhf)
    at (\yearx{2023}+1.2, \laneC*\laneheight) {Safe RLHF};
\node[paper, fill=thread3] (jailbreak)
    at (\yearx{2024}+1.2, \laneC*\laneheight) {Jailbreak};

% =====================================================================
% Thread 4: Reasoning & Verification (lane y=-3)
% =====================================================================
\def\laneD{-3}
\node[lanelabel, text=thread4] at (-0.3, \laneD*\laneheight)
    {T4: Reasoning \&\\Verification};

\node[paper, fill=thread4] (cot)
    at (\yearx{2022}+0.3, \laneD*\laneheight) {CoT};
\node[paper, fill=thread4] (prm)
    at (\yearx{2023}+0.5, \laneD*\laneheight) {PRM};
\node[paper, fill=thread4] (oone)
    at (\yearx{2024}+0.0, \laneD*\laneheight) {o1};
\node[paper, fill=thread4] (deepseekr1)
    at (\yearx{2025}-0.3, \laneD*\laneheight) {R1};
\node[paper, fill=thread4] (othree)
    at (\yearx{2025}+0.8, \laneD*\laneheight) {o3};
\node[paper, fill=thread4] (v32)
    at (\yearx{2026}+0.0, \laneD*\laneheight) {V3.2};

% =====================================================================
% Thread 5: Multimodal Post-training (lane y=-4)
% =====================================================================
\def\laneE{-4}
\node[lanelabel, text=thread5] at (-0.3, \laneE*\laneheight)
    {T5: Multimodal\\Post-training};

\node[paper, fill=thread5] (llava)
    at (\yearx{2023}-0.3, \laneE*\laneheight) {LLaVA};
\node[paper, fill=thread5] (vlmrlhf)
    at (\yearx{2024}+0.0, \laneE*\laneheight) {VLM-RLHF};
\node[paper, fill=thread5] (halluc)
    at (\yearx{2025}+0.3, \laneE*\laneheight) {Anti-Halluc.};

% =====================================================================
% Thread 6: Tool Use & Agentic (lane y=-5)
% =====================================================================
\def\laneF{-5}
\node[lanelabel, text=thread6] at (-0.3, \laneF*\laneheight)
    {T6: Tool Use \&\\Agentic};

\node[paper, fill=thread6] (toolformer)
    at (\yearx{2022}+1.5, \laneF*\laneheight) {Toolformer};
\node[paper, fill=thread6] (funccall)
    at (\yearx{2023}+1.5, \laneF*\laneheight) {Func.\ Call};
\node[paper, fill=thread6] (agenttune)
    at (\yearx{2024}+1.5, \laneF*\laneheight) {AgentTune};
\node[paper, fill=thread6] (mcp)
    at (\yearx{2025}+1.3, \laneF*\laneheight) {MCP/Agents};

% =====================================================================
% Thread 7: Efficiency & Data Engineering (lane y=-6)
% =====================================================================
\def\laneG{-6}
\node[lanelabel, text=thread7] at (-0.3, \laneG*\laneheight)
    {T7: Efficiency \&\\Data Engineering};

\node[paper, fill=thread7] (lora)
    at (\yearx{2021}+1.0, \laneG*\laneheight) {LoRA};
\node[paper, fill=thread7] (datasel)
    at (\yearx{2023}+0.8, \laneG*\laneheight) {Data Selection};
\node[paper, fill=thread7] (synthdata)
    at (\yearx{2024}+0.5, \laneG*\laneheight) {Synthetic Data};

% =====================================================================
% Cross-thread arrows (dashed, idea flow)
% =====================================================================

% T1 -> T2: InstructGPT-SFT feeds into InstructGPT-RLHF
\draw[crossarrow, thread1] (instructgpt-sft.south) -- (instructgpt-rlhf.north);

% T2 -> T3: RLHF ideas feed into Constitutional AI
\draw[crossarrow, thread2] (instructgpt-rlhf.south) -- (cai.north)
    node[arrowlabel, midway, right=1pt] {reward};

% T2 -> T4: Reward modeling feeds reasoning verification (PRM)
\draw[crossarrow, thread2] (dpo.south)
    to[out=-90, in=90] (prm.north);

% T4 -> T2: Reasoning verification inspires GRPO
\draw[crossarrow, thread4] (prm.east)
    to[out=0, in=180] (grpo.south west)
    node[arrowlabel, midway, above=1pt] {verification};

% T1 -> T5: SFT pipeline adopted by LLaVA
\draw[crossarrow, thread1] (selfinstruct.south)
    to[out=-90, in=90] (llava.north)
    node[arrowlabel, near start, right=1pt] {SFT pipeline};

% T2 -> T5: RLHF adopted for VLM alignment
\draw[crossarrow, thread2] (dpo.south)
    to[out=-100, in=90] (vlmrlhf.north)
    node[arrowlabel, near end, left=1pt] {preference};

% T1 -> T6: Instruction tuning enables tool use
\draw[crossarrow, thread1] (selfinstruct.south)
    to[out=-90, in=90] (toolformer.north);

% T7 -> T1: LoRA enables efficient SFT
\draw[crossarrow, thread7] (lora.north)
    to[out=90, in=-90] (flan.south)
    node[arrowlabel, near start, right=1pt] {LoRA};

% T7 -> T2: Efficient training enables broader RLHF
\draw[crossarrow, thread7] (datasel.north)
    to[out=90, in=-90] (grpo.south);

% T3 -> T4: Safety constraints shape reasoning alignment
\draw[crossarrow, thread3] (saferlhf.south)
    to[out=-90, in=90] (oone.north)
    node[arrowlabel, midway, right=1pt] {safety};

% T4 -> T6: Reasoning enables agent planning
\draw[crossarrow, thread4] (oone.south)
    to[out=-90, in=90] (agenttune.north)
    node[arrowlabel, midway, right=1pt] {reasoning};

% T2 -> T4: RLVR (DAPO) feeds reasoning training
\draw[crossarrow, thread2] (dapo.south)
    to[out=-90, in=90] (othree.north)
    node[arrowlabel, near start, right=1pt] {RLVR};

% T4 -> T5: Reasoning RL migrates to multimodal (R1-VL, MM-Eureka)
\draw[crossarrow, thread4] (deepseekr1.south)
    to[out=-90, in=90] (halluc.north east);

% T4 -> T6: Thinking models empower agent planning
\draw[crossarrow, thread4] (othree.south)
    to[out=-90, in=90] (mcp.north);

\end{tikzpicture}%
}% end resizebox
\caption{Post-training 领域的 Hybrid Problem-Evolution Graph 全景视图。
七条水平泳道分别对应七条演化线程，
x 轴为时间线（2017--2026），
节点为各线程内的 landmark 工作，
虚线箭头表示跨线程的思想流动与技术迁移。}
\label{fig:overview}
\end{figure}
